\scanpage{130}\mspage{75}
la première implique la seconde, et l'inverse est vrai si $X$ est propre sur $S$, et dans ce cas, \deleted{\unsure{les}}la condition équiva\deleted{\unsure{lent}}ut à ce que $X$ soit \deleted{\unsure{dans|des}} $\text{\uline{Fibhot}}\,S$. Et \unsure{enfin|en fait}, si $X$ est déjà dans $\text{\uline{Fib}}\ \text{\uline{hot}}\,S$, alors les deux conditions sont automatiquement satisfaites.)
% slanted
\margin{\unsure{On en est donc|On en vient donc} dans le diagramme \unsure{à} introduire aussi $\text{\uline{Liss}}_0 S,\ \text{\uline{Prop}}_0 S$ !}

% slanted
\margin{le diagramme des \uline{quinze} catégories \unsure{remarquables} !}
\diagram{\begin{tikzcd}[ampersand replacement=\&,row sep=2pt,column sep=1pt,font=\scriptsize]
 \text{\uline{Fibsc}}_0\,S\arrow[dd,hook]\arrow[dr]\&{}\&{}\&{}\&{}\&{}\&\text{\uline{Cofibsc}}_0\,S\arrow[dd,hook]\arrow[dl]\\
 {}\&\deleted{\unsure{F}}\text{\uline{Fib}}_0\,S\arrow[dd,hook]\arrow[dr,hook]\&{}\&\mathrm{Parf}\,S\arrow[dl,hook]\arrow[dd,hook]\arrow[dr,hook]\&{}\&\text{\uline{Cofib}}_0\,S\arrow[dl,hook]\arrow[dd,hook]\&{}\\
 \text{\uline{Fibsc}}\,S\arrow[dr,hook]\&{}\&\text{\uline{Liss}}_0\,S\arrow[dd,hook]\arrow[dr,hook]\&{}\&\text{\uline{Prop}}_0\,S\arrow[dd,hook]\arrow[dl,hook]\&{}\&\text{\uline{Cofibsc}}\,S\arrow[dl,hook]\\
 {}\&\text{\uline{Fib}}\,S\arrow[dr,hook]\&{}\&\text{\uline{Fibhot}}\,S\arrow[dd,hook]\arrow[uuulll,dashed,"\Phi_0" description]\arrow[uuurrr,dashed,"\Psi_0" description]\&{}\&\text{\uline{Cofib}}\,S\arrow[dl,hook]\&{}\\
 {}\&{}\&\text{\uline{Liss}}\,S\arrow[dr,hook]\&{}\&\text{\uline{Prop}}\,S\arrow[dl]\&{}\&{}\\
 {}\&{}\&{}\&\mathrm{Cat}_{/S}\arrow[uuulll,dashed,bend left=20,"\Phi"']\arrow[uuurrr,dashed,bend right=20,"\Psi"']\&{}\&{}\&{}.
\end{tikzcd}}

De plus, si deux éléments $\xi\in\mathrm{Ob}\,\mathrm{HOT}_W(S)$, $\xi'\in\mathrm{HOT}^\circ_W(S)$ proviennent d'un même $\unsure{y|\eta}\in\mathrm{Ob}\,\mathrm{HOT}^{\unsure{\mathrm{lc}|\ell\mathrm{c}}}_W(S)$, alors les deux foncteurs \deleted{\unsure{associés} \unsure{sur}} correspondants
\[
 (\Pi S)^\circ\xrightarrow{\varphi}\mathrm{Hot}_W\xleftarrow{\varphi'}\Pi S
\]
sont \uline{associés}, i.e. \deleted{chacun est} que chacun est isomorphe à ``l'inverse'' de l'autre.
\[
 \left\{\begin{aligned}
 \varphi^\circ(s)&=\varphi(s)\\
 \varphi^\circ(u)&=\varphi(u)^{-1},
 \end{aligned}\right.
\]

\textcircled{VI} (NB Dans la nouvelle version du diagramme, où j'ai rajouté des catégories $\text{\uline{Liss}}_0 S$ et $\text{\uline{Prop}}_0 S$ pour rendre le \unsure{truc|tracé} plus transparent, les catégories appelées ``bordantes'' deviennent le ``rez de chaussée'', et les catégories ``internes'', le ``premier étage'' ou ``étage''. Tous les \inserted{\unsure{7}} carrés du diagramme sont cartésiens. Toutes les flèches sont des inclusions pl. fid, sauf les \deleted{\unsure{carrés}} flèches horizontales \unsure{extérieures} des catégories de Cat-fibrés ou de Cat-cofibrés \uline{scindés}, avec morphismes exclusivement cartésiens ou cocartésiens et \unsure{mêmes|même} compatibles avec les scindages, vers les cat. corresp. \unsure{admettant|donnant} des structures non scindées)
% archive sheet number: 129
